% META: {"id": "TSPE-CONTIN-CO-013", "chapitre": "TSPE-CONTINUITE", "type_objet": "corrige", "exercice_ref": "TSPE-CONTIN-EX-013", "capacites_codes": ["C1"], "sources_inspiration": [], "status": "generated"}
% BEGIN-VERIFY
% from sympy import symbols, diff, Rational, nsolve, sqrt, exp, S
% x = symbols('x')
% f = 2*x**3 - 9*x**2 + 12*x
% assert (diff(f, x) - 6*(x-1)*(x-2)).expand() == 0
% assert f.subs(x, 2) == 4 and f.subs(x, 3) == 9
% assert diff(f, x).subs(x, 2) == 0
% assert all(diff(f, x).subs(x, v) > 0 for v in (Rational(21,10), Rational(5,2), 3))
% assert f.subs(x, Rational(26,10)) == Rational(5512,1000)
% assert f.subs(x, Rational(27,10)) == Rational(6156,1000)
% assert f.subs(x, Rational(26,10)) < 6 < f.subs(x, Rational(27,10))
% END-VERIFY
\begin{corrige}{TSPE-CONTIN-EX-013}

\textbf{Q1.} $f'(x) = 6x^2 - 18x + 12$. En developpant, $6(x-1)(x-2) = 6(x^2 - 3x + 2) = 6x^2 - 18x + 12$ : les deux expressions coincident.

\textbf{Q2.} Sur $[2\,;3]$ : $x - 1 > 0$ et $x - 2 \geq 0$, l'annulation n'ayant lieu qu'en $x = 2$. Donc $f'(x) > 0$ sur $\left]2\,;3\right]$ et $f'(2) = 0$. La derivee ne s'annulant qu'en une borne, $f$ est strictement croissante sur $[2\,;3]$.

\textbf{Q3.} $f(2) = 16 - 36 + 24 = 4$ et $f(3) = 54 - 81 + 36 = 9$.

\textbf{Q4.} $f$ est continue et strictement croissante sur $[2\,;3]$, et $6$ est compris entre $f(2) = 4$ et $f(3) = 9$ : l'equation $f(x) = 6$ admet une unique solution $\alpha$ dans $[2\,;3]$. De plus $f(2{,}6) = 5{,}512 < 6$ et $f(2{,}7) = 6{,}156 > 6$, donc $2{,}6 < \alpha < 2{,}7$.
\end{corrige}
